\section{What would change the logarithmic order?}
\label{sec:research}

The geometric menu removes a genuine bottleneck in the submitted
constant calculation. It does not remove the partition-growth mechanism
that produces the square root of the logarithm. It is useful to identify
the remaining task precisely, and to separate it from optimizations that
can only improve constants under the present interface.

\subsection{The alternative crossing interface in the public paper}
Anthropic's Appendix A gives a general crossing-piece construction with
a parameter $\rho$. A polynomial-size family of exact crossing pieces
whose row and column factorization bit counts sum to
\[
 4\rho\lambda+O(\log\lambda)
\]
gives logarithmic denominator power $\log_2(1/\rho)$. Appendix B realizes
$\rho=1/\sqrt2$ using a graded Heisenberg representation and ordinary
partition growth~\cite{anthropic-paper}. The full interface includes exact
arithmetic size, piece construction, entry evaluation, and an adaptive
choice of the cheaper row or column aggregation.

Appendix H already identifies the following research direction: a compatible
representation with
\[
 \log\dim V_E\asymp E^\alpha
\]
would suggest $\rho=2^{\alpha-1}$ and hence logarithmic denominator power
$1-\alpha$. This observation belongs to that source. A representation with
$\alpha<1/2$ would be a promising route to an order improvement \emph{if}
it also supplied all exact operator identities, locality, coefficient-size
bounds, and effective evaluation guarantees. A small graded dimension by
itself does not supply those properties.

The most useful research target after this review is therefore an exact,
efficiently evaluable two-sided crossing decomposition with
$\rho<1/\sqrt2$, or an equally explicit replacement for the Cauchy
weighted-mode interface. We do not construct either object here.

\subsection{Why merely keeping square-indexed oscillators fails}
The reviewed earlier second-round note already warned that assigning square
weights to the oscillator modes breaks the necessary degree closure
\cite{degree-barrier}. A different tempting change is to delete all
nonsquare modes. It does improve a bare dimension count, but has a simple
locality obstruction.

Let $p_\square(E)$ be the number of partitions of $E$ into square parts.
Then
\[
 \log p_\square(E)=\Theta(E^{1/3}).
\]
For a lower bound, choose arbitrary subsets of
$2^2,3^2,\ldots,J^2$ with $J=cE^{1/3}$ for a sufficiently small fixed
$c>0$, and fill the remaining weight with ones. For an upper bound, use
the generating function and the Gaussian sum:
\begin{align*}
 \log p_\square(E)
 &\le\theta E+\sum_{k\ge1}\frac1k\sum_{j\ge1}e^{-\theta k j^2}\\
 &\le\theta E+O(\theta^{-1/2});
\end{align*}
take $\theta=E^{-2/3}$. This numerical growth improvement is not enough.

\begin{proposition}[A scoped finite-pole obstruction]
Let $S$ be an infinite set of positive integers with arbitrarily long gaps,
and put $a_m=m$ when $|m|\in S$ and $a_m=0$ otherwise. The restricted
scalar Heisenberg-current commutator
\[
 C_S(z,w)=\sum_{m\in\mathbb Z}a_m z^{-m-1}w^{m-1}
\]
is not annihilated by $(z-w)^H$ for any finite nonnegative integer $H$.
\end{proposition}
\begin{proof}
If that product vanished, coefficient comparison would give
\[
 \sum_{j=0}^H(-1)^j\binom Hj a_{n-j}=0
 \quad\text{for every integer }n.
\]
For $H>0$, $H$ consecutive zero values would force all later values to be
zero by this recurrence. An arbitrarily long gap followed by another
element of $S$ is a contradiction. The case $H=0$ is immediate.
For square indices, there is the explicit witness
$n=(H+1)^2$: all preceding $H$ indices are nonsquares, and the displayed
coefficient equals $n\ne0$.
\end{proof}
This blocks the naive deletion of oscillator modes in the current used
by the source construction. It is not a lower bound against all vertex
algebras, all sparse gradings, all crossing decompositions, or all SAT
algorithms. The included finite witness program is a check on the formula;
the proposition itself is an all-size proof.

\subsection{Further constant improvements and their limit}
The support-averaging problem has an independent finite-geometric form.
One seeks menus for which the sums of all possible operation supports
have small Euclidean norm. The affine-plane construction works because
different parallel classes are transverse, and this controls every
intermediate support as well as the final triangle. It is reasonable to
search other designs or weighted mixtures of plans. However, the coefficient
we prove is a valid certificate for the supplied menu, not a proof that
the menu or the rational schedule is globally optimal.

A pointwise cheapest-plan bound can also be strictly better than the
uniform arithmetic average used by Lemma~\ref{lem:average}. Establishing
that improvement would require a new uniform bound on
$\min_j\max_{v\in\mathcal V_j}v\cdot g$ over the admissible nonnegative
entropy vectors. Numerical optimization without a rigorous certificate
would not justify a stronger constant. Likewise, letting the group count
grow with input size requires controlling the number of budget allocations,
all pair directions, packet packing, and table construction uniformly;
our fixed-$r$ proof does not do so.

The twelve-group menu was found by finite search and then certified
independently. Its validity depends only on the frozen plans and the
integer certificate. Replaying the search is not necessary to verify the
result. No negative search result in this report is promoted to an
impossibility theorem.
